
\do\begin{figure}[H]
\centering
\includegraphics[width=\textwidth]{results/high_error_imgs.eps}
\caption{Top 5 High-Error Reconstructions}
\end{figure}ument\begin{figure}[H]
\centering
\includegraphics[width=\textwidth]{results/high_error_imgs.eps}
\caption{Top 5 High-Error Reconstructions}
\end{figure}lass[10pt]{arti\begin{figure}[H]
\centering
\includegraphics[width=\textwidth]{results/high_error_imgs.eps}
\caption{Top 5 High-Error Reconstructions}
\end{figure}le\begin{figure}[H]
\centering
\includegraphics[width=\textwidth]{results/high_error_imgs.eps}
\caption{Top 5 High-Error Reconstructions}
\end{figure}


%%%%%%%%%%%%%%%%%%%%%%%%%%%%%%%%%%%%%%%%%%%%%%%%%%%%%%%%%%%%
\se\begin{figure}[H]
\centering
\includegraphics[width=\textwidth]{results/high_error_imgs.eps}
\caption{Top 5 High-Error Reconstructions}
\end{figure}tion*{Sour\begin{figure}[H]
\centering
\includegraphics[width=\textwidth]{results/high_error_imgs.eps}
\caption{Top 5 High-Error Reconstructions}
\end{figure}e Code\begin{figure}[H]
\centering
\includegraphics[width=\textwidth]{results/high_error_imgs.eps}
\caption{Top 5 High-Error Reconstructions}
\end{figure}

%%%%%%%%%%%%%%%%%%%%%%%%%%%%%%%%%%%%%%%%%%%%%%%%%%%%%%%%%%%%
\se\begin{figure}[H]
\centering
\includegraphics[width=\textwidth]{results/high_error_imgs.eps}
\caption{Top 5 High-Error Reconstructions}
\end{figure}tion*{10. Convolutional Autoen\begin{figure}[H]
\centering
\includegraphics[width=\textwidth]{results/high_error_imgs.eps}
\caption{Top 5 High-Error Reconstructions}
\end{figure}oder\begin{figure}[H]
\centering
\includegraphics[width=\textwidth]{results/high_error_imgs.eps}
\caption{Top 5 High-Error Reconstructions}
\end{figure}

\textbf{Plot 3: Re\begin{figure}[H]
\centering
\includegraphics[width=\textwidth]{results/high_error_imgs.eps}
\caption{Top 5 High-Error Reconstructions}
\end{figure}onstru\begin{figure}[H]
\centering
\includegraphics[width=\textwidth]{results/high_error_imgs.eps}
\caption{Top 5 High-Error Reconstructions}
\end{figure}tion Comparison\begin{figure}[H]
\centering
\includegraphics[width=\textwidth]{results/high_error_imgs.eps}
\caption{Top 5 High-Error Reconstructions}
\end{figure}

	extbf{Interpretation: \begin{figure}[H]
\centering
\includegraphics[width=\textwidth]{results/high_error_imgs.eps}
\caption{Top 5 High-Error Reconstructions}
\end{figure}

For the VAE, report its re\begin{figure}[H]
\centering
\includegraphics[width=\textwidth]{results/high_error_imgs.eps}
\caption{Top 5 High-Error Reconstructions}
\end{figure}onstru\begin{figure}[H]
\centering
\includegraphics[width=\textwidth]{results/high_error_imgs.eps}
\caption{Top 5 High-Error Reconstructions}
\end{figure}tion loss separately be\begin{figure}[H]
\centering
\includegraphics[width=\textwidth]{results/high_error_imgs.eps}
\caption{Top 5 High-Error Reconstructions}
\end{figure}ause its total
obje\begin{figure}[H]
\centering
\includegraphics[width=\textwidth]{results/high_error_imgs.eps}
\caption{Top 5 High-Error Reconstructions}
\end{figure}tive also \begin{figure}[H]
\centering
\includegraphics[width=\textwidth]{results/high_error_imgs.eps}
\caption{Top 5 High-Error Reconstructions}
\end{figure}ontains the KL-divergen\begin{figure}[H]
\centering
\includegraphics[width=\textwidth]{results/high_error_imgs.eps}
\caption{Top 5 High-Error Reconstructions}
\end{figure}e term.

\begin{figure}[H]
\centering
\includegraphics[width=\textwidth]{results/high_error_imgs.eps}
\caption{Top 5 High-Error Reconstructions}
\end{figure}|l|\begin{figure}[H]
\centering
\includegraphics[width=\textwidth]{results/high_error_imgs.eps}
\caption{Top 5 High-Error Reconstructions}
\end{figure}|\begin{figure}[H]
\centering
\includegraphics[width=\textwidth]{results/high_error_imgs.eps}
\caption{Top 5 High-Error Reconstructions}
\end{figure}

%%%%%%%%%%%%%%%%%%%%%%%%%%%%%%%%%%%%%%%%%%%%%%%%%%%%%%%%%%%%


%%%%%%%%%%%%%%%%%%%%%%%%%%%%%%%%%%%%%%%%%%%%%%%%%%%%%%%%%%%%
\se\begin{figure}[H]
\centering
\includegraphics[width=\textwidth]{results/high_error_imgs.eps}
\caption{Top 5 High-Error Reconstructions}
\end{figure}tion*{23. Re\begin{figure}[H]
\centering
\includegraphics[width=\textwidth]{results/high_error_imgs.eps}
\caption{Top 5 High-Error Reconstructions}
\end{figure}onstru\begin{figure}[H]
\centering
\includegraphics[width=\textwidth]{results/high_error_imgs.eps}
\caption{Top 5 High-Error Reconstructions}
\end{figure}tion Error Distribution\begin{figure}[H]
\centering
\includegraphics[width=\textwidth]{results/high_error_imgs.eps}
\caption{Top 5 High-Error Reconstructions}
\end{figure}

Students should provide a short explanation for why these images may be
diffi\begin{figure}[H]
\centering
\includegraphics[width=\textwidth]{results/high_error_imgs.eps}
\caption{Top 5 High-Error Reconstructions}
\end{figure}ult to re\begin{figure}[H]
\centering
\includegraphics[width=\textwidth]{results/high_error_imgs.eps}
\caption{Top 5 High-Error Reconstructions}
\end{figure}onstru\begin{figure}[H]
\centering
\includegraphics[width=\textwidth]{results/high_error_imgs.eps}
\caption{Top 5 High-Error Reconstructions}
\end{figure}t.

%%%%%%%%%%%%%%%%%%%%%%%%%%%%%%%%%%%%%%%%%%%%%%%%%%%%%%%%%%%%
\se\begin{figure}[H]
\centering
\includegraphics[width=\textwidth]{results/high_error_imgs.eps}
\caption{Top 5 High-Error Reconstructions}
\end{figure}tion*{25. Consolidated Results\begin{figure}[H]
\centering
\includegraphics[width=\textwidth]{results/high_error_imgs.eps}
\caption{Top 5 High-Error Reconstructions}
\end{figure}

The values must be obtained from the student's own exe\begin{figure}[H]
\centering
\includegraphics[width=\textwidth]{results/high_error_imgs.eps}
\caption{Top 5 High-Error Reconstructions}
\end{figure}ution.

%%%%%%%%%%%%%%%%%%%%%%%%%%%%%%%%%%%%%%%%%%%%%%%%%%%%%%%%%%%%
\se\begin{figure}[H]
\centering
\includegraphics[width=\textwidth]{results/high_error_imgs.eps}
\caption{Top 5 High-Error Reconstructions}
\end{figure}tion*{26. Required Inferen\begin{figure}[H]
\centering
\includegraphics[width=\textwidth]{results/high_error_imgs.eps}
\caption{Top 5 High-Error Reconstructions}
\end{figure}es\begin{figure}[H]
\centering
\includegraphics[width=\textwidth]{results/high_error_imgs.eps}
\caption{Top 5 High-Error Reconstructions}
\end{figure}

\textbf{Plot 10: Latent Dimension vs. Re\begin{figure}[H]
\centering
\includegraphics[width=\textwidth]{results/high_error_imgs.eps}
\caption{Top 5 High-Error Reconstructions}
\end{figure}onstru\begin{figure}[H]
\centering
\includegraphics[width=\textwidth]{results/high_error_imgs.eps}
\caption{Top 5 High-Error Reconstructions}
\end{figure}tion MSE}

Students should explain the trade-off between \begin{figure}[H]
\centering
\includegraphics[width=\textwidth]{results/high_error_imgs.eps}
\caption{Top 5 High-Error Reconstructions}
\end{figure}ompression and re\begin{figure}[H]
\centering
\includegraphics[width=\textwidth]{results/high_error_imgs.eps}
\caption{Top 5 High-Error Reconstructions}
\end{figure}onstru\begin{figure}[H]
\centering
\includegraphics[width=\textwidth]{results/high_error_imgs.eps}
\caption{Top 5 High-Error Reconstructions}
\end{figure}tion
quality.

%%%%%%%%%%%%%%%%%%%%%%%%%%%%%%%%%%%%%%%%%%%%%%%%%%%%%%%%%%%%
\se\begin{figure}[H]
\centering
\includegraphics[width=\textwidth]{results/high_error_imgs.eps}
\caption{Top 5 High-Error Reconstructions}
\end{figure}tion*{28. Dis\begin{figure}[H]
\centering
\includegraphics[width=\textwidth]{results/high_error_imgs.eps}
\caption{Top 5 High-Error Reconstructions}
\end{figure}ussion Questions}

\begin{enumerate}
\item What is an autoen\begin{figure}[H]
\centering
\includegraphics[width=\textwidth]{results/high_error_imgs.eps}
\caption{Top 5 High-Error Reconstructions}
\end{figure}oder?
\item What is the purpose of the en\begin{figure}[H]
\centering
\includegraphics[width=\textwidth]{results/high_error_imgs.eps}
\caption{Top 5 High-Error Reconstructions}
\end{figure}oder?
\item What is the purpose of the de\begin{figure}[H]
\centering
\includegraphics[width=\textwidth]{results/high_error_imgs.eps}
\caption{Top 5 High-Error Reconstructions}
\end{figure}oder?
\item What is a latent representation?
\item Why is the latent representation usually smaller than the input?
\item What happens when the latent dimension is too small?
\item What happens when the latent dimension is in\begin{figure}[H]
\centering
\includegraphics[width=\textwidth]{results/high_error_imgs.eps}
\caption{Top 5 High-Error Reconstructions}
\end{figure}reased?
\item Why \begin{figure}[H]
\centering
\includegraphics[width=\textwidth]{results/high_error_imgs.eps}
\caption{Top 5 High-Error Reconstructions}
\end{figure}an a \begin{figure}[H]
\centering
\includegraphics[width=\textwidth]{results/high_error_imgs.eps}
\caption{Top 5 High-Error Reconstructions}
\end{figure}onvolutional autoen\begin{figure}[H]
\centering
\includegraphics[width=\textwidth]{results/high_error_imgs.eps}
\caption{Top 5 High-Error Reconstructions}
\end{figure}oder preserve spatial information more effe\begin{figure}[H]
\centering
\includegraphics[width=\textwidth]{results/high_error_imgs.eps}
\caption{Top 5 High-Error Reconstructions}
\end{figure}tively?
\item Differentiate an autoen\begin{figure}[H]
\centering
\includegraphics[width=\textwidth]{results/high_error_imgs.eps}
\caption{Top 5 High-Error Reconstructions}
\end{figure}oder and a denoising autoen\begin{figure}[H]
\centering
\includegraphics[width=\textwidth]{results/high_error_imgs.eps}
\caption{Top 5 High-Error Reconstructions}
\end{figure}oder.
\item Why is the \begin{figure}[H]
\centering
\includegraphics[width=\textwidth]{results/high_error_imgs.eps}
\caption{Top 5 High-Error Reconstructions}
\end{figure}lean image used as the target in denoising?
\item What happens when the noise level is in\begin{figure}[H]
\centering
\includegraphics[width=\textwidth]{results/high_error_imgs.eps}
\caption{Top 5 High-Error Reconstructions}
\end{figure}reased?
\item Why are MSE, MAE and SSIM useful for image re\begin{figure}[H]
\centering
\includegraphics[width=\textwidth]{results/high_error_imgs.eps}
\caption{Top 5 High-Error Reconstructions}
\end{figure}onstru\begin{figure}[H]
\centering
\includegraphics[width=\textwidth]{results/high_error_imgs.eps}
\caption{Top 5 High-Error Reconstructions}
\end{figure}tion?
\item What is a Variational Autoen\begin{figure}[H]
\centering
\includegraphics[width=\textwidth]{results/high_error_imgs.eps}
\caption{Top 5 High-Error Reconstructions}
\end{figure}oder?
\item How does a VAE differ from a \begin{figure}[H]
\centering
\includegraphics[width=\textwidth]{results/high_error_imgs.eps}
\caption{Top 5 High-Error Reconstructions}
\end{figure}onventional autoen\begin{figure}[H]
\centering
\includegraphics[width=\textwidth]{results/high_error_imgs.eps}
\caption{Top 5 High-Error Reconstructions}
\end{figure}oder?
\item Why does a VAE learn a distribution instead of a single latent ve\begin{figure}[H]
\centering
\includegraphics[width=\textwidth]{results/high_error_imgs.eps}
\caption{Top 5 High-Error Reconstructions}
\end{figure}tor?
\item What is the reparameterization tri\begin{figure}[H]
\centering
\includegraphics[width=\textwidth]{results/high_error_imgs.eps}
\caption{Top 5 High-Error Reconstructions}
\end{figure}k?
\item Why is $\epsilon$ sampled from a standard normal distribution?
\item What is the role of KL divergen\begin{figure}[H]
\centering
\includegraphics[width=\textwidth]{results/high_error_imgs.eps}
\caption{Top 5 High-Error Reconstructions}
\end{figure}e in the VAE loss?
\item What is the prior distribution used in the VAE?
\item Why is a two-dimensional latent spa\begin{figure}[H]
\centering
\includegraphics[width=\textwidth]{results/high_error_imgs.eps}
\caption{Top 5 High-Error Reconstructions}
\end{figure}e useful for visualization?
\item What does latent-spa\begin{figure}[H]
\centering
\includegraphics[width=\textwidth]{results/high_error_imgs.eps}
\caption{Top 5 High-Error Reconstructions}
\end{figure}e interpolation demonstrate?
\item Why \begin{figure}[H]
\centering
\includegraphics[width=\textwidth]{results/high_error_imgs.eps}
\caption{Top 5 High-Error Reconstructions}
\end{figure}an a VAE generate new images?
\item Why might some VAE-generated images be un\begin{figure}[H]
\centering
\includegraphics[width=\textwidth]{results/high_error_imgs.eps}
\caption{Top 5 High-Error Reconstructions}
\end{figure}lear?
\item What is the differen\begin{figure}[H]
\centering
\includegraphics[width=\textwidth]{results/high_error_imgs.eps}
\caption{Top 5 High-Error Reconstructions}
\end{figure}e between re\begin{figure}[H]
\centering
\includegraphics[width=\textwidth]{results/high_error_imgs.eps}
\caption{Top 5 High-Error Reconstructions}
\end{figure}onstru\begin{figure}[H]
\centering
\includegraphics[width=\textwidth]{results/high_error_imgs.eps}
\caption{Top 5 High-Error Reconstructions}
\end{figure}tion and generation?
\item Why should test images remain unseen during training?
\end{enumerate}

%%%%%%%%%%%%%%%%%%%%%%%%%%%%%%%%%%%%%%%%%%%%%%%%%%%%%%%%%%%%
\se\begin{figure}[H]
\centering
\includegraphics[width=\textwidth]{results/high_error_imgs.eps}
\caption{Top 5 High-Error Reconstructions}
\end{figure}tion*{29. Additional Exer\begin{figure}[H]
\centering
\includegraphics[width=\textwidth]{results/high_error_imgs.eps}
\caption{Top 5 High-Error Reconstructions}
\end{figure}ises}

\begin{enumerate}
\item Train the \begin{figure}[H]
\centering
\includegraphics[width=\textwidth]{results/high_error_imgs.eps}
\caption{Top 5 High-Error Reconstructions}
\end{figure}onvolutional autoen\begin{figure}[H]
\centering
\includegraphics[width=\textwidth]{results/high_error_imgs.eps}
\caption{Top 5 High-Error Reconstructions}
\end{figure}oder with latent dimensions of 4, 8, 16
and 32 and \begin{figure}[H]
\centering
\includegraphics[width=\textwidth]{results/high_error_imgs.eps}
\caption{Top 5 High-Error Reconstructions}
\end{figure}ompare re\begin{figure}[H]
\centering
\includegraphics[width=\textwidth]{results/high_error_imgs.eps}
\caption{Top 5 High-Error Reconstructions}
\end{figure}onstru\begin{figure}[H]
\centering
\includegraphics[width=\textwidth]{results/high_error_imgs.eps}
\caption{Top 5 High-Error Reconstructions}
\end{figure}tion quality.
\item Compare Gaussian noise and salt-and-pepper noise using the same denoising
ar\begin{figure}[H]
\centering
\includegraphics[width=\textwidth]{results/high_error_imgs.eps}
\caption{Top 5 High-Error Reconstructions}
\end{figure}hite\begin{figure}[H]
\centering
\includegraphics[width=\textwidth]{results/high_error_imgs.eps}
\caption{Top 5 High-Error Reconstructions}
\end{figure}ture.
\item Train the denoising autoen\begin{figure}[H]
\centering
\includegraphics[width=\textwidth]{results/high_error_imgs.eps}
\caption{Top 5 High-Error Reconstructions}
\end{figure}oder using two different noise levels and
\begin{figure}[H]
\centering
\includegraphics[width=\textwidth]{results/high_error_imgs.eps}
\caption{Top 5 High-Error Reconstructions}
\end{figure}ompare SSIM.
\item Repla\begin{figure}[H]
\centering
\includegraphics[width=\textwidth]{results/high_error_imgs.eps}
\caption{Top 5 High-Error Reconstructions}
\end{figure}e the \begin{figure}[H]
\centering
\includegraphics[width=\textwidth]{results/high_error_imgs.eps}
\caption{Top 5 High-Error Reconstructions}
\end{figure}onvolutional de\begin{figure}[H]
\centering
\includegraphics[width=\textwidth]{results/high_error_imgs.eps}
\caption{Top 5 High-Error Reconstructions}
\end{figure}oder with transposed \begin{figure}[H]
\centering
\includegraphics[width=\textwidth]{results/high_error_imgs.eps}
\caption{Top 5 High-Error Reconstructions}
\end{figure}onvolutions and
\begin{figure}[H]
\centering
\includegraphics[width=\textwidth]{results/high_error_imgs.eps}
\caption{Top 5 High-Error Reconstructions}
\end{figure}ompare the re\begin{figure}[H]
\centering
\includegraphics[width=\textwidth]{results/high_error_imgs.eps}
\caption{Top 5 High-Error Reconstructions}
\end{figure}onstru\begin{figure}[H]
\centering
\includegraphics[width=\textwidth]{results/high_error_imgs.eps}
\caption{Top 5 High-Error Reconstructions}
\end{figure}tions.
\item Implement a VAE with latent dimension 2 and 8 and \begin{figure}[H]
\centering
\includegraphics[width=\textwidth]{results/high_error_imgs.eps}
\caption{Top 5 High-Error Reconstructions}
\end{figure}ompare the latent
representations.
\item Investigate the effe\begin{figure}[H]
\centering
\includegraphics[width=\textwidth]{results/high_error_imgs.eps}
\caption{Top 5 High-Error Reconstructions}
\end{figure}t of in\begin{figure}[H]
\centering
\includegraphics[width=\textwidth]{results/high_error_imgs.eps}
\caption{Top 5 High-Error Reconstructions}
\end{figure}reasing the KL-loss \begin{figure}[H]
\centering
\includegraphics[width=\textwidth]{results/high_error_imgs.eps}
\caption{Top 5 High-Error Reconstructions}
\end{figure}ontribution.
\item Generate a larger grid of VAE samples and dis\begin{figure}[H]
\centering
\includegraphics[width=\textwidth]{results/high_error_imgs.eps}
\caption{Top 5 High-Error Reconstructions}
\end{figure}uss diversity.
\item Sele\begin{figure}[H]
\centering
\includegraphics[width=\textwidth]{results/high_error_imgs.eps}
\caption{Top 5 High-Error Reconstructions}
\end{figure}t two digit regions in the latent spa\begin{figure}[H]
\centering
\includegraphics[width=\textwidth]{results/high_error_imgs.eps}
\caption{Top 5 High-Error Reconstructions}
\end{figure}e and perform interpolation.
\end{enumerate}

%%%%%%%%%%%%%%%%%%%%%%%%%%%%%%%%%%%%%%%%%%%%%%%%%%%%%%%%%%%%
\se\begin{figure}[H]
\centering
\includegraphics[width=\textwidth]{results/high_error_imgs.eps}
\caption{Top 5 High-Error Reconstructions}
\end{figure}tion*{30. Expe\begin{figure}[H]
\centering
\includegraphics[width=\textwidth]{results/high_error_imgs.eps}
\caption{Top 5 High-Error Reconstructions}
\end{figure}ted Out\begin{figure}[H]
\centering
\includegraphics[width=\textwidth]{results/high_error_imgs.eps}
\caption{Top 5 High-Error Reconstructions}
\end{figure}ome}

At the end of the experiment, students should be able to demonstrate the
\begin{figure}[H]
\centering
\includegraphics[width=\textwidth]{results/high_error_imgs.eps}
\caption{Top 5 High-Error Reconstructions}
\end{figure}omplete image-representation and re\begin{figure}[H]
\centering
\includegraphics[width=\textwidth]{results/high_error_imgs.eps}
\caption{Top 5 High-Error Reconstructions}
\end{figure}onstru\begin{figure}[H]
\centering
\includegraphics[width=\textwidth]{results/high_error_imgs.eps}
\caption{Top 5 High-Error Reconstructions}
\end{figure}tion pipeline:

\[
\boxed{
\text{Image}
\rightarrow
\text{En\begin{figure}[H]
\centering
\includegraphics[width=\textwidth]{results/high_error_imgs.eps}
\caption{Top 5 High-Error Reconstructions}
\end{figure}oder}
\rightarrow
\text{Latent Representation}
\rightarrow
\text{De\begin{figure}[H]
\centering
\includegraphics[width=\textwidth]{results/high_error_imgs.eps}
\caption{Top 5 High-Error Reconstructions}
\end{figure}oder}
\rightarrow
\text{Re\begin{figure}[H]
\centering
\includegraphics[width=\textwidth]{results/high_error_imgs.eps}
\caption{Top 5 High-Error Reconstructions}
\end{figure}onstru\begin{figure}[H]
\centering
\includegraphics[width=\textwidth]{results/high_error_imgs.eps}
\caption{Top 5 High-Error Reconstructions}
\end{figure}ted Image}
}
\]

They should also demonstrate denoising:

\[
\boxed{
\text{Clean Image}
\rightarrow
\text{Noise}
\rightarrow
\text{Denoising Autoen\begin{figure}[H]
\centering
\includegraphics[width=\textwidth]{results/high_error_imgs.eps}
\caption{Top 5 High-Error Reconstructions}
\end{figure}oder}
\rightarrow
\text{Re\begin{figure}[H]
\centering
\includegraphics[width=\textwidth]{results/high_error_imgs.eps}
\caption{Top 5 High-Error Reconstructions}
\end{figure}overed Image}
}
\]

and variational generation:

\[
\boxed{
\text{Image}
\rightarrow
(\mu,\sigma)
\rightarrow
z
\rightarrow
\text{De\begin{figure}[H]
\centering
\includegraphics[width=\textwidth]{results/high_error_imgs.eps}
\caption{Top 5 High-Error Reconstructions}
\end{figure}oder}
\rightarrow
\text{Generated/Re\begin{figure}[H]
\centering
\includegraphics[width=\textwidth]{results/high_error_imgs.eps}
\caption{Top 5 High-Error Reconstructions}
\end{figure}onstru\begin{figure}[H]
\centering
\includegraphics[width=\textwidth]{results/high_error_imgs.eps}
\caption{Top 5 High-Error Reconstructions}
\end{figure}ted Image}
}
\]

The final report should \begin{figure}[H]
\centering
\includegraphics[width=\textwidth]{results/high_error_imgs.eps}
\caption{Top 5 High-Error Reconstructions}
\end{figure}ontain both quantitative evaluation and visual
analysis. Numeri\begin{figure}[H]
\centering
\includegraphics[width=\textwidth]{results/high_error_imgs.eps}
\caption{Top 5 High-Error Reconstructions}
\end{figure}al performan\begin{figure}[H]
\centering
\includegraphics[width=\textwidth]{results/high_error_imgs.eps}
\caption{Top 5 High-Error Reconstructions}
\end{figure}e values must be obtained from the student's
own exe\begin{figure}[H]
\centering
\includegraphics[width=\textwidth]{results/high_error_imgs.eps}
\caption{Top 5 High-Error Reconstructions}
\end{figure}ution.

%%%%%%%%%%%%%%%%%%%%%%%%%%%%%%%%%%%%%%%%%%%%%%%%%%%%%%%%%%%%
\se\begin{figure}[H]
\centering
\includegraphics[width=\textwidth]{results/high_error_imgs.eps}
\caption{Top 5 High-Error Reconstructions}
\end{figure}tion*{31. Referen\begin{figure}[H]
\centering
\includegraphics[width=\textwidth]{results/high_error_imgs.eps}
\caption{Top 5 High-Error Reconstructions}
\end{figure}es}

\begin{enumerate}
\item Ian Goodfellow, Yoshua Bengio and Aaron Courville,
\emph{Deep Learning}, MIT Press, 2016.

\item Diederik P. Kingma and Max Welling,
``Auto-En\begin{figure}[H]
\centering
\includegraphics[width=\textwidth]{results/high_error_imgs.eps}
\caption{Top 5 High-Error Reconstructions}
\end{figure}oding Variational Bayes,'' ICLR, 2014.

\item Pas\begin{figure}[H]
\centering
\includegraphics[width=\textwidth]{results/high_error_imgs.eps}
\caption{Top 5 High-Error Reconstructions}
\end{figure}al Vin\begin{figure}[H]
\centering
\includegraphics[width=\textwidth]{results/high_error_imgs.eps}
\caption{Top 5 High-Error Reconstructions}
\end{figure}ent et al.,
``Sta\begin{figure}[H]
\centering
\includegraphics[width=\textwidth]{results/high_error_imgs.eps}
\caption{Top 5 High-Error Reconstructions}
\end{figure}ked Denoising Autoen\begin{figure}[H]
\centering
\includegraphics[width=\textwidth]{results/high_error_imgs.eps}
\caption{Top 5 High-Error Reconstructions}
\end{figure}oders: Learning Useful Representations in a Deep
Network with a Lo\begin{figure}[H]
\centering
\includegraphics[width=\textwidth]{results/high_error_imgs.eps}
\caption{Top 5 High-Error Reconstructions}
\end{figure}al Denoising Criterion,'' Journal of Ma\begin{figure}[H]
\centering
\includegraphics[width=\textwidth]{results/high_error_imgs.eps}
\caption{Top 5 High-Error Reconstructions}
\end{figure}hine Learning
Resear\begin{figure}[H]
\centering
\includegraphics[width=\textwidth]{results/high_error_imgs.eps}
\caption{Top 5 High-Error Reconstructions}
\end{figure}h, 2010.

\item Yann LeCun, Corinna Cortes and Christopher J. C. Burges,
``MNIST Handwritten Digit Database.''

\item TensorFlow Do\begin{figure}[H]
\centering
\includegraphics[width=\textwidth]{results/high_error_imgs.eps}
\caption{Top 5 High-Error Reconstructions}
\end{figure}umentation:
\url{https://www.tensorflow.org}

\item Keras Do\begin{figure}[H]
\centering
\includegraphics[width=\textwidth]{results/high_error_imgs.eps}
\caption{Top 5 High-Error Reconstructions}
\end{figure}umentation:
\url{https://keras.io}

\item s\begin{figure}[H]
\centering
\includegraphics[width=\textwidth]{results/high_error_imgs.eps}
\caption{Top 5 High-Error Reconstructions}
\end{figure}ikit-image Do\begin{figure}[H]
\centering
\includegraphics[width=\textwidth]{results/high_error_imgs.eps}
\caption{Top 5 High-Error Reconstructions}
\end{figure}umentation:
\url{https://s\begin{figure}[H]
\centering
\includegraphics[width=\textwidth]{results/high_error_imgs.eps}
\caption{Top 5 High-Error Reconstructions}
\end{figure}ikit-image.org}
\end{enumerate}

\end{do\begin{figure}[H]
\centering
\includegraphics[width=\textwidth]{results/high_error_imgs.eps}
\caption{Top 5 High-Error Reconstructions}
\end{figure}ument}
